\documentclass{article}
\usepackage[utf8]{inputenc}
\usepackage[spanish]{babel}
\usepackage{amsmath, amssymb, amsthm}

\title{Resumen de Definiciones: El Significado de la Recursión}
\author{}
\date{}

\newtheorem{definition}{Definición}

\begin{document}

\maketitle

\section{Teoría de Órdenes Parciales}

\begin{definition}[Orden Parcial / Poset]
Un conjunto $X$ junto con una relación $\le$ es un \textit{poset} si la relación es reflexiva, antisimétrica y transitiva.
\end{definition}

\begin{definition}[Orden Discreto]
Es el orden parcial definido por la relación de igualdad ($x \le y \iff x = y$).
\end{definition}

\begin{definition}[Lifting]
Dado un poset $X$, su lifting $X_{\perp}$ se define como $X \cup \{\perp\}$, donde el nuevo orden satisface $x \le y$ si $x = \perp$ o $x \le_X y$.
\end{definition}

\begin{definition}[Orden Llano]
Se obtiene al aplicar el \textit{lifting} a un conjunto con orden discreto.
\end{definition}

\begin{definition}[Orden en Espacio de Funciones]
Dadas funciones $f, g: X \to Y$, se define $f \le g$ si y solo si $\forall x \in X, f(x) \le_Y g(x)$.
\end{definition}

\section{Espacios de Funciones y sus Propiedades de Orden}

\begin{definition}[Espacio Ordenado de Funciones]
Si $(Y, \le_Y)$ es un orden parcial y $X$ es un conjunto, el espacio de funciones de $X \to Y$ es un orden parcial donde la relación de orden entre dos funciones $f, g \in X \to Y$ está definida de manera puntual:
\begin{equation}
    f \le g \iff \forall x \in X, f(x) \le_Y g(x)
\end{equation}
.
\end{definition}

\begin{definition}[Preservación de Estructura]
El espacio de funciones hereda las propiedades de estructura del codominio:
\begin{itemize}
    \item Si $Y$ es un \textbf{predominio}, entonces el espacio de funciones $X \to Y$ también es un predominio.
    \item Si $D$ es un \textbf{dominio}, entonces $X \to D$ también lo es. En este caso, el elemento mínimo del espacio de funciones es la función constante que devuelve $\perp$ para toda entrada.
\end{itemize}
\end{definition}

\begin{definition}[Funciones Parciales y Orden]
Las funciones parciales de $X \rightharpoonup Y$ pueden ordenarse mediante la contención de sus gráficos (vistas como subconjuntos de $P(X \times Y)$), lo cual es equivalente a decir que $f \le g$ si siempre que $f(x)$ está definido, entonces $g(x)$ también lo está y $f(x) = g(x)$. Esta noción se formaliza mediante la totalización de funciones en el espacio $X \to Y_{\perp}$.
\end{definition}

\section{Cadenas y Supremos}

\begin{definition}[Supremo]
El supremo de un subconjunto $Q \subseteq P$ es el elemento $sup(Q) \in P$ que es cota superior de $Q$ y es menor o igual que cualquier otra cota superior de $Q$.
\end{definition}

\begin{definition}[Cadena]
Una secuencia infinita de elementos $p_0 \le p_1 \le p_2 \le \dots$. 
\begin{itemize}
    \item \textbf{Interesante}: Si el conjunto de sus elementos es infinito.
    \item \textbf{No interesante}: Si el conjunto es finito (se vuelve constante a partir de un punto).
\end{itemize}
\end{definition}

\section{Predominios y Dominios}

\begin{definition}[Predominio]
Un orden parcial $P$ es un predominio si toda cadena en $P$ tiene supremo. Equivalentemente, si toda cadena interesante tiene supremo.
\end{definition}

\begin{definition}[Dominio]
Es un predominio que posee un elemento mínimo (generalmente denotado como $\perp$).
\end{definition}

\section{Funciones y Morfismos}

\begin{definition}[Monotonía]
Una función $f: P \to Q$ es monótona si preserva el orden: $x \le_P y \implies f(x) \le_Q f(y)$.
\end{definition}

\begin{definition}[Continuidad]
Una función es continua si es monótona y preserva el supremo de las cadenas interesantes.
\end{definition}

\begin{definition}[Función Estricta]
Una función $f$ es estricta si preserva el elemento mínimo, es decir, $f(\perp) = \perp$.
\end{definition}

\section{Puntos Fijos}

\begin{definition}[Teorema del Menor Punto Fijo]
Sea $D$ un dominio y $F: D \to D$ una función continua. Entonces existe el menor punto fijo de $F$, el cual es el supremo de la cadena $\perp \le F(\perp) \le F^2(\perp) \le \dots$.
\end{definition}

\end{document}